% Additive incorporation for Article I; the preceding source is unchanged.
% RESULTADO_RECUPERADO: integral, chapter 22, extensive source
% colaboracion/partes_i_ii/source/base_residencias/
% capitulo_22_c20_lineas_2613_3270.tex, lines 176--576.
% Equivalence 2/3: public/residencias/capitulo_22_voa_leech_moonshine.tex,
% lines 55--127. Order two and Fricke: hmt/15d_voa_moonshine_orden_dos_rev11.tex.
% Prerequisites resident in I: code C_W, matrix A_W, N(A_2^12),
% origin o_*, vector v_alpha of norm54, N_alpha recognized as Leech,
% lattice VOA and FLM presentation in preceding sections.
% New bibliography: chenlamshimakura2016, abelamyamada2017,
% kmconwaynorton1979.

\subsection{Ternary orientation of the Leech neighbour}
\label{km:orden-tres-reticular}

The order-two presentation constructed above uses negation
of the Leech lattice. The same lattice, obtained from the incidence marked
by \(K\) and the dodecaphasic support of \(\alpha\), admits a second
presentation compatible with the ternary character of the construction.
To make this continuity explicit, the gluing is retained, its order-three
isometry is constructed, and preservation of the already selected neighbour
is verified. This step requires no alteration of \(K\), its rational reading,
or the preceding definition of \(\alpha\). The construction presents the lattice operations and the precise conditions of the classical theorems used.

We work in the basis of simple roots of each \(A_2\) factor, with
\begin{equation}
 B_2=\begin{pmatrix}2&-1\\-1&2\end{pmatrix},\qquad
 \mathsf C=\begin{pmatrix}0&-1\\1&-1\end{pmatrix},\qquad
 \mathsf R=\begin{pmatrix}0&1\\1&0\end{pmatrix},\qquad
 \rho=\begin{pmatrix}1\\1\end{pmatrix}.
 \label{km:coxeter-reflexion}
\end{equation}
Multiplication gives
\[
 \mathsf C^3=I_2,\quad
 \mathsf C^2+\mathsf C+I_2=0,\quad
 \mathsf C^{\mathsf T}B_2\mathsf C=B_2,\quad
 \mathsf R\mathsf C\mathsf R=\mathsf C^{-1},\quad
 \mathsf R^{\mathsf T}B_2\mathsf R=B_2.
\]
These identities fix both the metric and the orientation of the action.
In the discriminant \(A_2^*/A_2\cong\mathbb F_3\), we take
the class of \(\varpi=(2,1)^{\mathsf T}/3\) as generator. We have
\(\mathsf C\varpi-\varpi=(-1,0)^{\mathsf T}\), so
\(\mathsf C\) acts trivially on the discriminant; the reflection
\(\mathsf R\) reverses its sign.

The Paley--Witt matrix already constructed produces a full-support
element by the following operation:
\begin{equation}
 \begin{gathered}
 q_*=(1,1,1,1,1,1),\qquad q_*A_W=(2,1,1,1,1,1),\\
 c_*=(q_*,q_*A_W)
 =(1,1,1,1,1,1,2,1,1,1,1,1)\in C_W.
 \end{gathered}
 \label{km:palabra-total}
\end{equation}
Every coordinate of \(c_*\) is nonzero. For the twelve positions,
define the frames and the isometry
\begin{equation}
 P_b(c_*)=
 \begin{cases}\mathsf R,&(c_*)_b=1,\\I_2,&(c_*)_b=2,\end{cases}
 \qquad
 g_c=\bigoplus_{b=1}^{12}P_b(c_*)\mathsf C P_b(c_*)^{-1}.
 \label{km:isometria-ternaria}
\end{equation}
The symbol \(P_b(c_*)\) here denotes an invertible frame of an
\(A_2\) factor, not the three-dimensional isotypic projector \(P_3\).

\begin{proposition}[Order-three isometry on the marked neighbour]
\label{km:estabilidad-vecino}
The operator \(g_c\) is an order-three isometry with no nonzero
fixed vectors. It preserves the Niemeier lattice \(N=N(C_W)\) and the
neighbour \(N_\alpha\) constructed in \eqref{exc:vecino}.
\end{proposition}
\begin{proof}
Each block of \(g_c\) is \(\mathsf C\) or \(\mathsf C^{-1}\);
its minimal polynomial is \(X^2+X+1\). This proves its order and
the absence of fixed vectors. Both blocks act trivially on the
discriminant, so they preserve every gluing class and therefore
the lattice \(N\).

Write \(v=v_\alpha=4\rho_{o_*}+\sum_{b\ne o_*}\rho_b\),
with the frames and incidence origin already fixed. Its radial
coefficients are \(a_{o_*}=4\) and \(a_b=1\) for \(b\ne o_*\).
All satisfy \(a_b\equiv1\pmod3\). In one factor,
\[
 (\mathsf C-I_2)\rho=(-2,-1)^{\mathsf T},\qquad
 (\mathsf C^{-1}-I_2)\rho=(-1,-2)^{\mathsf T}.
\]
After division by three, their discriminant classes are
\(2\) and \(1\), respectively. Consequently, the vectors
\begin{equation}
 y=\frac{g_cv-v}{3},\qquad z=\frac{g_c^{-1}v-v}{3}
 \label{km:desplazamientos}
\end{equation}
have discriminant words \(c_*\) and \(-c_*\), both in
\(C_W\). Hence \(y,z\in N\). The local product is
\[
 \left\langle
 \frac{(\mathsf C^{\pm1}-I_2)a_b\rho}{3},a_b\rho
 \right\rangle=-a_b^2,
\]
and summing gives
\(\langle y,v\rangle=\langle z,v\rangle=-(16+11)=-27\).
Thus, \(y,z\in N_0=\{x\in N:\langle x,v\rangle\equiv0\pmod3\}\).
For \(x\in N_0\),
\[
 \langle g_cx,v\rangle
 =\langle x,g_c^{-1}v\rangle
 =\langle x,v\rangle+3\langle x,z\rangle\equiv0\pmod3.
\]
The argument applied to \(g_c^{-1}\) proves \(g_cN_0=N_0\).
Finally, \(g_c(v/3)=v/3+y\) preserves the class generating
\(N_\alpha=N_0+\mathbb Z(v/3)\). We obtain
\(g_cN_\alpha=N_\alpha\).
\end{proof}

The proof preserves more than the lattice dimension: it uses the
gluing code, its orientation, the flag origin, and the congruence of the
radial vector. These data explain why the second presentation
remains linked to the dodecaphasic record from which the construction began.

\subsection{Ternary trace, twisted weight, and order-three extension}
\label{km:orbifold-tres}

Let \(\Lambda=N_\alpha\), and let \(\widehat g_c\) be the standard
lift of \(g_c\) to the lattice algebra
\(V_\Lambda=M(1)\otimes\mathbb C_\varepsilon[\Lambda]\), with the
central extension and cocycle already specified. The odd order allows
this lift to be chosen of order three. Each complex \(A_2\) factor
contributes the eigenvalues \(\zeta_3\) and \(\zeta_3^2\), where
\(\zeta_3^2+\zeta_3+1=0\).

\begin{proposition}[Graded trace and type of the lift]
\label{km:traza-tipo}
For \(j=1,2\),
\begin{equation}
 \operatorname{Tr}_{V_\Lambda}
       \bigl(\widehat g_c^{\,j}q^{L_0-1}\bigr)
 =t_3(\tau)
 :=q^{-1}\prod_{n\ge1}(1+q^n+q^{2n})^{-12}.
 \label{km:tres-traza}
\end{equation}
The lowest weight of the nontrivial twisted modules is \(4/3\), and
the lift is of type zero.
\end{proposition}
\begin{proof}
Only the zero vector contributes to the lattice part of the trace,
because \(g_c\) and \(g_c^2\) have no nonzero fixed vectors.
At each oscillator degree, the trace of symmetric powers is the inverse
of the determinant \(I-g_c^jq^n\). Each \(A_2\) factor contributes
\(1+q^n+q^{2n}\), proving \eqref{km:tres-traza}.

The twisted-vacuum weight formula for an order-three lattice
automorphism, applied to the two twelve-dimensional eigenspaces, gives
\[
 \rho_{\widehat g_c}
 =\frac{1}{4\cdot3^2}
   \bigl(1\cdot2\cdot12+2\cdot1\cdot12\bigr)=\frac43.
\]
The type is the class of \(3^2\rho_{\widehat g_c}=12\) modulo
three, which is zero. The weight formula is a result of twisted-module
theory; the multiplicities and isometry supplied to it
have been constructed here.
\end{proof}

\begin{theorem}[Second presentation of the Moonshine algebra]
\label{km:moonshine-tres}
The type-zero cyclic extension
\begin{equation}
 V_{(3)}=\operatorname{Orb}_{\mathbb Z/3}
                 (V_\Lambda,\widehat g_c)
 \label{km:extension-orden3}
\end{equation}
is isomorphic to the vertex operator algebra \(V^\natural\).
\end{theorem}
\begin{proof}
Theorem \ref{exc:teorema-leech} provides a positive,
even, unimodular, rootless lattice of rank twenty-four.
Proposition~\ref{km:estabilidad-vecino} constructs an order-three
fixed-point-free isometry on that same lattice;
Proposition~\ref{km:traza-tipo} verifies the lift and its type.
The Chen--Lam--Shimakura theorem on the order-three orbifold
of the Leech lattice algebra then applies
\cite{chenlamshimakura2016}. This identification is also covered
by the uniform treatment of Abe--Lam--Yamada
\cite{abelamyamada2017}. This application follows construction
of the lattice data; equality of characters does not replace the
hypotheses of the theorem used.
\end{proof}

To specify the dual symmetry, write
\(V_{(3)}=W_0\oplus W_1\oplus W_2\), where \(W_i\) is the
integral sector selected in the \(\widehat g_c^{\,i}\)-twisted module.
The product satisfies
\(Y(W_i,z)W_j\subset W_{i+j\bmod3}((z))\). Hence
\begin{equation}
 g^\vee|_{W_i}=\zeta_3^i I
 \label{km:dual-ciclico}
\end{equation}
defines an order-three automorphism. Indeed, the factor on a
product is \(\zeta_3^{i+j}=\zeta_3^i\zeta_3^j\), and the vacuum and
conformal vector belong to \(W_0\). The classical identification of
this automorphism in the Monster is class \(3B\), with series
\(T_{3B}=t_3+12\) \cite{kmconwaynorton1979,borcherds}.
This normalization can be checked in the construction. The lattice
character is \(\operatorname{ch}V_\Lambda=J+24\): its initial pole
is \(q^{-1}\), its weight-one space has dimension 24, and the
weight-zero modular character is determined by these data.
If \(F\) is the character of the fixed subspace and \(H\) that of each
integral twisted sector, projection onto invariants and conjugation
of the two sectors give
\[
 F=\frac{J+24+2t_3}{3},\qquad F+2H=J,
 \qquad \operatorname{Tr}g^\vee q^{L_0-1}=F-H=t_3+12.
\]
Its domain is the extended algebra, whereas
\(\widehat g_c\) acts on the preceding lattice algebra.

\subsection{Equivalence of the order-two and order-three presentations}
\label{km:equivalencia-dos-tres}

Denote by
\[
 V_{(2)}=(V_\Lambda)^+\oplus(V_\Lambda^T)^+
\]
the FLM presentation in \eqref{exc:flm}. Both presentations preserve
the same Leech neighbour as their antecedent. Their extensions employ
different automorphisms and vertex products determined by
the respective lifting data.

\begin{theorem}[Graded equivalence and transport of automorphisms]
\label{km:equivalencia-voa}
Given vertex operator algebra isomorphisms
\[
 \phi_2:V_{(2)}\xrightarrow{\ \cong\ }V^\natural,
 \qquad
 \phi_3:V_{(3)}\xrightarrow{\ \cong\ }V^\natural,
\]
the composition
\begin{equation}
 \Phi_{23}=\phi_2^{-1}\phi_3:
 V_{(3)}\xrightarrow{\ \cong\ }V_{(2)}
 \label{km:Phi23}
\end{equation}
preserves the vacuum, conformal vector, grading, and vertex
product. Conjugation by \(\Phi_{23}\) identifies the automorphism
groups and preserves their graded traces.
\end{theorem}
\begin{proof}
The composition of the two isomorphisms preserves the stated
structures. In particular,
\[
 \Phi_{23}Y_{(3)}(a,z)b
 =Y_{(2)}(\Phi_{23}a,z)\Phi_{23}b,
 \qquad
 \Phi_{23}L_0^{(3)}=L_0^{(2)}\Phi_{23}.
\]
If \(h\) is an automorphism of \(V_{(3)}\), then
\(h' =\Phi_{23}h\Phi_{23}^{-1}\) is an automorphism of
\(V_{(2)}\). Equality of traces on each finite-dimensional homogeneous
component gives
\[
 \operatorname{Tr}_{V_{(3)}}h q^{L_0-1}
 =\operatorname{Tr}_{V_{(2)}}h' q^{L_0-1}.
\]
\end{proof}

The choice of \(\phi_2\) and \(\phi_3\) forms part of the statement;
\(\Phi_{23}\) is not presented as a canonical identification
independent of these choices. Nor does it transform an order-three
automorphism into the order-two involution: conjugation preserves
order. Its content is that two different extensions realize the
same graded structure, with explicit transport of their products,
representations, and traces.

\subsection{Eta functions, Fricke involution, and vacuum normalization}
\label{km:fricke}

For \(\operatorname{Im}\tau>0\), the traces of the two lattice
automorphisms admit the expressions
\begin{equation}
 \eta(\tau)=q^{1/24}\prod_{n\ge1}(1-q^n),\qquad
 t_2=\left(\frac{\eta(\tau)}{\eta(2\tau)}\right)^{24},\qquad
 t_3=\left(\frac{\eta(\tau)}{\eta(3\tau)}\right)^{12}.
 \label{km:eta-trazas}
\end{equation}
They follow directly from the factorizations
\(1-q^{2n}=(1-q^n)(1+q^n)\) and
\(1-q^{3n}=(1-q^n)(1+q^n+q^{2n})\). The order in \(q\) of the
level-three quotient, before raising it to the twelfth power, is
\[
 \frac1{24}-\frac3{24}=-\frac1{12};
 \qquad 12\left(-\frac1{12}\right)=-1.
\]
Here \(-1/12\) denotes the exact difference between the initial
exponents of two eta functions. The operation is specified and
does not require interpreting a divergent sum as a convergent one.
Agreement with other normalized defects in the corpus preserves
their respective realizations; it does not identify their operators
through equality of the scalar.

\begin{proposition}[Fricke involution on the order-two trace]
\label{km:fricke-prop}
Let \(w_2\tau=-1/(2\tau)\). Then
\begin{equation}
 t_2(w_2\tau)=\frac{2^{12}}{t_2(\tau)},\qquad
 T_{2A}=t_2+24+\frac{2^{12}}{t_2}
 \label{km:fricke-t2}
\end{equation}
and \(T_{2A}(w_2\tau)=T_{2A}(\tau)\).
\end{proposition}
\begin{proof}
The transformation law
\(\eta(-1/\tau)=(-i\tau)^{1/2}\eta(\tau)\) gives
\[
 \frac{\eta(-1/(2\tau))}{\eta(-1/\tau)}
 =\sqrt2\,\frac{\eta(2\tau)}{\eta(\tau)}.
\]
The twenty-fourth power yields the first identity. Substitution
interchanges the two variable terms of \(T_{2A}\) and fixes
the constant term. Recognition of this expression as the
McKay--Thompson series of class \(2A\) corresponds to the
level-two symmetrization of Conway--Norton
\cite[\S\S4 and 6, Table~3]{kmconwaynorton1979} and the
Moonshine theorem \cite{borcherds}.
\end{proof}

The modular uniformizations of levels two and three, with the
preceding normalizations, are
\begin{align}
 j&=\frac{(t_2+256)^3}{t_2^2},
 \label{km:j-nivel2}\\
 J=j-744
 &=t_3+12+\frac{10\cdot3^9}{t_3}
       +\frac{4\cdot3^{14}}{t_3^2}
       +\frac{3^{18}}{t_3^3}.
 \label{km:j-nivel3}
\end{align}
They are used as classical modular identities on the traces already
constructed, not as rules for selecting the code or the neighbour.
The second can also be written
\(j=(t_3+27)(t_3+243)^3/t_3^3\).
The direct product expansion of \eqref{km:tres-traza} begins with
\[
 t_3=q^{-1}-12+54q-76q^2-243q^3+1188q^4+O(q^5).
\]
In particular, \([q]J=54+10\cdot3^9=196884\); this is a
consequence of the traces and uniformization. The construction of
\(V^\natural\) retains its lattice and orbifold proof,
independently of this check of the first coefficient.

% BEGIN R27_MG09
\subsubsection{Ternary congruence of the trace and optimality of the exponent}
\label{mg09:congruencia}

The level-three identity permits comparison of every coefficient of
the constructed trace with the corresponding coefficient of its modular
realization. The appropriate domain for this consequence is the ring
of formal power series: each coefficient depends on finitely many
factors of the product, while the statement retains all degrees.

\begin{lemma}[Integral unit of the order-three product]
The series
\begin{equation}
 u(q)=q\,t_3(q)=\prod_{n\geq1}(1+q^n+q^{2n})^{-12}
 \label{mg09:unidad}
\end{equation}
belongs to \(1+q\mathbb Z[[q]]\), and its inverse belongs to the
same set.
\end{lemma}

\begin{proof}
Each polynomial \(1+q^n+q^{2n}\) has constant term one. For any
series \(v(q)=1+\sum_{r\geq1}v_rq^r\) with integer
coefficients, the equation \(v(q)w(q)=1\) recursively determines
\[
 w_0=1,\qquad w_r=-\sum_{i=1}^{r}v_iw_{r-i}\quad(r\geq1).
\]
This recurrence proves existence, uniqueness and integrality of the
inverse. Applied to each factor, it also proves integrality of its
inverse twelfth power. Moreover,
\((1+q^n+q^{2n})^{-12}\in1+q^n\mathbb Z[[q]]\).
To compute a coefficient of degree at most \(r\), only factors
with \(n\leq r\) are needed; the partial products stabilize
coefficientwise. Consequently, \eqref{mg09:unidad} defines an
integral series with constant term one. Applying the same recurrence
to \(u\) yields \(u^{-1}\in1+q\mathbb Z[[q]]\).
\end{proof}

\begin{proposition}[Exact uniform divisibility]
For the trace \(t_3\) and the series \(J=j-744\) related by
\eqref{km:j-nivel3},
\begin{equation}
 J-(t_3+12)\in3^9q\mathbb Z[[q]].
 \label{mg09:congruencia-exacta}
\end{equation}
The uniform exponent of divisibility by three is exactly nine:
\begin{equation}
 [q]\bigl(J-t_3-12\bigr)=10\cdot3^9=196830.
 \label{mg09:optimalidad}
\end{equation}
\end{proposition}

\begin{proof}
Since \(t_3=q^{-1}u\), the lemma gives
\(t_3^{-r}=q^ru^{-r}\in q^r\mathbb Z[[q]]\) for
\(r=1,2,3\). Identity \eqref{km:j-nivel3} becomes
\begin{equation}
 J-t_3-12
 =3^9q\left(10u^{-1}+4\cdot3^5q\,u^{-2}
                         +3^9q^2u^{-3}\right).
 \label{mg09:factorizacion}
\end{equation}
Every term in parentheses is an integral series, proving
\eqref{mg09:congruencia-exacta}. Its constant term is \(10\),
since \(u(0)=u^{-1}(0)=1\), and the last two terms contain
\(q\) and \(q^2\). This gives \eqref{mg09:optimalidad}. As
\(3\nmid10\), the coefficient of \(q\) has \(3\)-adic valuation
nine, establishing optimality of the uniform exponent.
\end{proof}

In particular, \([q^r]J\equiv[q^r]t_3\pmod{3^9}\) for every
\(r\geq1\). The quantifier follows from the formal identity and
the integral inversion proved in the lemma. Verification through any
finite degree is a specialization of this relation. The trace continues
to arise from the order-three automorphism of the lattice constructed
from HMT incidence; uniformization subsequently acts on that trace.
The congruence preserves this order and determines an exact relation
between its two graded readings.

% END R27_MG09

The three operations occurring in this extension have
different domains. The dual symmetry \(g^\vee\) acts on sectors
of the vertex algebra; Fricke acts on the modular parameter and
its function; radius duality interchanges momentum and winding
in the compact lattice sector. Their reciprocity formulas can
be coordinated through the corresponding maps, but the name
“duality” does not replace those maps. This distinction allows the
construction to continue towards dimensional screens while preserving
incidence, grading, and the origin of each operation.

% New bibitems appear in the article bibliography.
% Primary verification: arXiv 1606.05961v2 (Theorem 1.1);
% arXiv 1705.09022v4 (Sections 2--4);
% Conway--Norton, pp. 311 and 314--315, Table 3.
